\section{An explicit auxiliary congruence}\label{sec:lean}
One of the conditions on an auxiliary prime in Alpöge's source is a
simultaneous congruence modulo $16$ and an odd integer $M$ \cite{source}.
It has the following explicit representative.
\begin{prop}\label{prop:crt}
For positive odd $M$ and any nonnegative integer $d$,
\[
 d\equiv7\pmod {16},\quad M\mid d+1
 \quad\Longleftrightarrow\quad d\equiv8M-1\pmod{16M}.
\]
Moreover $\gcd(8M-1,16M)=1$ for every positive integer $M$.
\end{prop}
\begin{proof}
Write $d+1=Mk$. Since $M$ is odd it acts invertibly modulo $16$
and fixes the residue $8$. Hence $Mk\equiv8\pmod{16}$ is
equivalent to $k\equiv8\pmod{16}$. This yields the claimed class,
and reversing the substitution proves the converse. Every common
divisor of $8M-1$ and $16M$ divides $16M-2(8M-1)=2$, while $8M-1$
is odd.
\end{proof}
For example $M=33$ gives $d\equiv263\pmod{528}$. With $M=3$ one
obtains $d\equiv23\pmod{48}$, which selects every other integer on
the $23$rd spoke. A prime $q>3$ with $q\equiv3\pmod8$ must lie on
the $11$th or the $19$th spoke. These congruences prescribe candidate
classes. An estimate for primes in a prescribed interval with varying
$M$ requires additional analytic input.

The accompanying Lean file proves the two spoke statements, the
explicit congruence, coprimality, and existence of the residue class.
Sections~\ref{sec:local}--\ref{sec:explicit} have mathematical proofs
in the text. The proof and verification records of Alpöge and
Sutherland concern the separate upstream project \cite{alpoge,sutherland}.
